\newpage
\section{问题重述}

\subsection{问题背景}
近年来，随着\textbf{[领域背景]}快速发展，\textbf{[具体问题]}日益突出。根据\textbf{[数据来源]}统计，\textbf{[关键数据与趋势]}。如何科学有效地解决这一问题，已成为\textbf{[相关领域]}亟待攻克的难题。

\subsection{问题提出}
本文针对\textbf{[问题主题]}问题，提出以下三个子问题：

\textbf{问题一}: 基于\textbf{[数据类型]}，分析\textbf{[规律特征]}，建立\textbf{[模型类型1]}模型，实现\textbf{[目标1]}。设观测序列为$\{\y_t\}\_{t=1}^T$，需提取其\textbf{[特征模式]}。

\textbf{问题二}: 在问题一的基础上，构建\textbf{[模型类型2]}，综合考虑\textbf{[影响因素]}，优化\textbf{[决策变量]}$x \in \mathcal{X}$，实现\textbf{[目标2]}：
$$\min_{x \in \mathcal{X}} f(x)\quad\text{s.t.}\quad g_j(x) \leq 0, \h_k(x) = 0$$

\textbf{问题三}: 建立\textbf{[模型类型3]}，以\textbf{[目标函数]}为核心，考虑\textbf{[约束条件]}，求解\textbf{[最优解]}，为\textbf{[决策者]}提供科学依据。

\subsection{研究意义}
本研究对于\textbf{[理论意义]}和\textbf{[实际应用价值]}具有重要的理论与现实意义。
